\documentclass[a4paper,12pt]{article}

\usepackage{graphicx}
\usepackage{caption}
\usepackage{subcaption}
\usepackage{tikz}
\usepackage{pgf}
\usepackage{amsmath}
\usepackage{amssymb}
\usetikzlibrary{arrows.meta}
\usepackage[utf8]{inputenc}
\usepackage[english,greek]{babel}
\usepackage[hidelinks]{hyperref}
\usepackage{multirow}
\usepackage[normalem]{ulem}

\begin{document}
\selectlanguage{greek}

\begin{center}
    {\Large ΕΡΓΑΣΙΑ ΣΤΟ ΜΑΘΗΜΑ ΒΑΣΕΙΣ ΔΕΔΟΜΕΝΩΝ}\\[3cm]

    {\Large \selectlanguage{english}\textbf{CryptofolioDB}\selectlanguage{greek}}\\[0.5cm]

    {\large Μία προσωπική εφαρμογή επενδύσεων και παρακολούθησης κρυπτονομισμάτων }\\[3cm]

    {\Large Πρώτο Παραδοτέο}\\[3cm]

    {\large Ομάδα 39}\\[1cm]

    \begin{tabular}{l l l}
        Δασκαλόπουλος Αριστείδης & 10640 & \selectlanguage{english}\texttt{arisdask@ece.auth.gr}\selectlanguage{greek} \\
        Λαβδάκης Γεώργιος & 10672 & \selectlanguage{english}\texttt{glavdakis@ece.auth.gr}\selectlanguage{greek} \\
        Ρουσομάνης Γεώργιος & 10703 & \selectlanguage{english}\texttt{rousoman@ece.auth.gr}\selectlanguage{greek} \\
    \end{tabular}\\[3cm]

    {\large 19/11/2025}
\end{center}

\newpage

\tableofcontents

\newpage

\section{Εισαγωγή}

\subsection{Σκοπός Εφαρμογής}
Σκοπός της εφαρμογής είναι να παρέχει μία ολοκληρωμένη, ασφαλή και επεκτάσιμη πλατφόρμα για τη
διαχείριση ψηφιακών πορτοφολιών και επενδύσεων σε κρυπτονομίσματα. Το σύστημα επιτρέπει στους
χρήστες να οργανώνουν τα περιουσιακά τους στοιχεία, να παρακολουθούν την αξία των
κρυπτονομισμάτων σε πραγματικό χρόνο, να εκτελούν συναλλαγές αγοράς και πώλησης, καθώς και να
συνεργάζονται με άλλους χρήστες ως συνδικαιούχοι πορτοφολιών. Παράλληλα, προσφέρει εργαλεία
ειδοποιήσεων και αυτοματοποιημένης ενημέρωσης για σημαντικές μεταβολές στην αγορά.

\subsection{Περιγραφή Εφαρμογής}
Η εφαρμογή λειτουργεί ως μια \selectlanguage{english}web-based\selectlanguage{greek} πλατφόρμα όπου κάθε εγγεγραμμένος χρήστης διαθέτει
τουλάχιστον ένα ψηφιακό πορτοφόλι και μπορεί να συνδέσει πολλαπλούς τραπεζικούς λογαριασμούς για
τη χρηματοδότηση των συναλλαγών του. Μέσω του γραφικού περιβάλλοντος, ο χρήστης έχει τη
δυνατότητα να:

\begin{itemize}
    \item δημιουργεί και να διαχειρίζεται πολλά πορτοφόλια
    \item παρακολουθεί τα διαθέσιμα υπόλοιπα και την αξία κάθε κρυπτονομίσματος
    \item εκτελεί αγορά ή πώληση κρυπτονομισμάτων μέσω εξωτερικών \selectlanguage{english}API\selectlanguage{greek} ανταλλακτηρίων
    \item πραγματοποιεί μεταφορές χρημάτων από συνδεδεμένους τραπεζικούς λογαριασμούς
    \item ορίζει ειδοποιήσεις για συγκεκριμένα τιμολογιακά όρια
    \item διαχειρίζεται την πρόσβαση άλλων χρηστών στο πορτοφόλι του μέσω προσκλήσεων
\end{itemize}

Η εφαρμογή απαιτεί την αποθήκευση δεδομένων που σχετίζονται με:

\begin{itemize}
    \item διαπιστευτήρια χρηστών και στοιχεία ταυτοποίησης
    \item ψηφιακά πορτοφόλια και συσχετίσεις μεταξύ χρηστών και πορτοφολιών
    \item τραπεζικούς λογαριασμούς και μεταφορές χρημάτων
    \item κρυπτονομίσματα, χαρακτηριστικά τους και είδη συναλλαγών
    \item ιστορικά διακυμάνσεων τιμών σε μορφή \selectlanguage{english}candlestick\selectlanguage{greek}
    \item ειδοποιήσεις τιμών και το αντίστοιχο καθεστώς ενεργοποίησης
\end{itemize}

\subsection{Απαιτήσεις Εφαρμογής σε Δεδομένα}

\begin{itemize}

  \item \textbf{Χρήστες:} 
  Αναμένεται ότι το σύστημα θα εξυπηρετεί από 1.000 έως 10.000 εγγεγραμμένους χρήστες 
  στην αρχική φάση λειτουργίας.

  \item \textbf{Πορτοφόλια ανά χρήστη:} 
  Κατά μέσο όρο 2--3 πορτοφόλια ανά χρήστη. 
  Συνολικά: 2.000 -- 30.000 εγγραφές.

  \item \textbf{\selectlanguage{english}Crypto balances\selectlanguage{greek}:} 
  Κατά μέσο όρο 5 διαφορετικά κρυπτονομίσματα ανά πορτοφόλι.  
  Συνεπώς αναμένονται 10.000 -- 150.000 εγγραφές.

  \item \textbf{Ιστορικό τιμών (\selectlanguage{english}CryptoHistory\selectlanguage{greek}):}  
  Για κάθε κρυπτονόμισμα αποθηκεύονται περίπου:
  \begin{itemize}
    \item 1.440 εγγραφές την ημέρα (ένα λεπτό ανά μέτρηση)
    \item Για 50 κρυπτονομίσματα: περίπου 72.000 εγγραφές/ημέρα
    \item Σε ετήσια βάση: πάνω από 26 εκατομμύρια εγγραφές
  \end{itemize}

  \item \textbf{\selectlanguage{english}Candlesticks\selectlanguage{greek}:}  
  Με \selectlanguage{english}interval\selectlanguage{greek} 1 ώρας για 50 νομίσματα: 
  50 \(\times\) 24 = 1200 εγγραφές/ημέρα

  \item \textbf{Συναλλαγές Κρυπτονομισμάτων:}  
  Σε εφαρμογές \selectlanguage{english}crypto-wallet\selectlanguage{greek} καταγράφονται τυπικά:
  \begin{itemize}
    \item 5--20 συναλλαγές ανά ενεργό χρήστη την ημέρα
    \item Για 2.000 ενεργούς χρήστες: 10.000 -- 40.000 εγγραφές/ημέρα
  \end{itemize}

  \item \textbf{Τραπεζικές Συναλλαγές:}  
  ~1--2 μεταφορές ανά χρήστη την εβδομάδα.
  Συνολικά: περίπου 10.000--20.000 εγγραφές/έτος.

\end{itemize}

\newpage

\section{Κατηγορίες Χρηστών και Απαιτήσεις τους}

\subsection*{Διαχειριστής Συστήματος (\selectlanguage{english}Administrator)}
Ο διαχειριστής συστήματος έχει πλήρη δικαιώματα στη βάση δεδομένων και ασχολείται με κινήσεις υψηλής σημασίας. Δεν χρησιμοποιεί την εφαρμογή ως απλός χρήστης αλλά, επεμβαίνει όταν χρειάζονται ενέργειες συντήρησης, όπως:
\begin{itemize}
\item δημιουργία και διαγραφή λογαριασμών χρηστών
\item εκτέλεση ελέγχων ασφαλείας
\item δημιουργία αντιγράφων ασφαλείας της βάσης
\item αποκατάσταση ή διόρθωση δεδομένων όπου απαιτείται
\end{itemize}

\subsection*{Διαχειριστής Πορτοφολιού}
Ο διαχειριστής πορτοφολιού έχει τον πλήρη έλεγχο ενός πορτοφολιού. Οι απαιτήσεις
του περιλαμβάνουν:
\begin{itemize}
    \item δημιουργία νέου πορτοφολιού
    \item πρόσκληση άλλων χρηστών ως συνδικαιούχων
    \item ανάκληση δικαιωμάτων ή περιορισμό πρόσβασης
    \item εκτέλεση συναλλαγών αγοράς και πώλησης
    \item προβολή όλων των τραπεζικών μεταφορών και υπολοίπων
\end{itemize}

\subsection*{Συνδικαιούχος Πορτοφολιού}
Ο συνδικαιούχος έχει περιορισμένα αλλά λειτουργικά δικαιώματα εντός του πορτοφολιού:
\begin{itemize}
    \item εκτέλεση συναλλαγών στο πορτοφόλι
    \item προβολή διαθέσιμων υπολοίπων και ιστορικού συναλλαγών
    \item αδυναμία διαχείρισης άλλων συνδικαιούχων
\end{itemize}

\newpage

\section{Μοντέλο Οντοτήτων/Συσχετίσεων}

\subsection{Γενική Περιγραφή}
Το υποκείμενο μοντέλο δεδομένων βασίζεται σε οντότητες που αναπαριστούν χρήστες, πορτοφόλια,
τραπεζικούς λογαριασμούς, κρυπτονομίσματα, συναλλαγές, ιστορικά τιμών και ειδοποιήσεις. Κεντρική
οντότητα αποτελεί ο \selectlanguage{english}\textbf{User}\selectlanguage{greek}, καθώς συνδέεται άμεσα 
ή έμμεσα με κάθε λειτουργικό μοντέλο της εφαρμογής.

Κάθε χρήστης μπορεί να κατέχει πολλά πορτοφόλια
(\selectlanguage{english}\textbf{Wallets}\selectlanguage{greek}) ενώ, κάθε πορτοφόλι μπορεί να έχει πολλούς
χρήστες με διαφορετικό ρόλο, διαχειριστές ή συνδικαιούχους και τον διαχειριστή συστήματος. Επιπλέον, οι χρήστες μπορούν να
συνδέουν πολλούς τραπεζικούς λογαριασμούς (\selectlanguage{english}\textbf{BankAccounts}\selectlanguage{greek}),
από τους οποίους εκτελούνται τραπεζικές συναλλαγές προς τα πορτοφόλια.

Ένα \selectlanguage{english}\textbf{Wallet}\selectlanguage{greek} περιέχει πολλαπλά \selectlanguage{english}\textbf{CryptoBalances}\selectlanguage{greek}, καθένα εκ των οποίων αφορά ένα
συγκεκριμένο \selectlanguage{english}\textbf{Cryptocurrency}\selectlanguage{greek}. Οι συναλλαγές συνδέονται πάντα με ένα 
\selectlanguage{english}CryptoBalance\selectlanguage{greek} και περιγράφουν την αγορά ή πώληση ενός κρυπτονομίσματος σε συγκεκριμένη τιμή και
χρονική στιγμή.

Το σύστημα καταγράφει αναλυτικά ιστορικά δεδομένα τιμών μέσω της οντότητας
\selectlanguage{english}\textbf{CandleStick}\selectlanguage{greek}, η οποία αποθηκεύει \selectlanguage{english}open/high/low/close\selectlanguage{greek} τιμές για κάθε κρυπτονόμισμα ανά
χρονική περίοδο. Παράλληλα, οι χρήστες μπορούν να ορίζουν όρια ενεργοποίησης
που βασίζονται στις τιμές των κρυπτονομισμάτων.

\subsubsection*{Υποθέσεις}
\begin{itemize}
    \item Δεν μπορεί ένας χρήστης να δημιουργήσει δύο διαφορετικούς λογαριασμούς με το ίδιο 
    \selectlanguage{english}email\selectlanguage{greek}.
    \item Το σύμβολο του κάθε κρυπτονομίσματος είναι μοναδικό.
    \item Ένα πορτοφόλι πρέπει να έχει \textbf{ακριβώς} έναν διαχειριστή.
    \item Κατά τη δημιουργία ενός νέου πορτοφολιού, ο χρήστης ορίζεται αυτόματα ως 
    \textit{διαχειριστής} του πορτοφολιού.
    \item Ένα νέο πορτοφόλι είναι αρχικά κενό. Για να πραγματοποιηθούν συναλλαγές κρυπτονομισμάτων, 
    απαιτείται πρώτα η μεταφορά χρημάτων σε αυτό από έναν συνδεδεμένο τραπεζικό λογαριασμό.
    \item Κάθε συναλλαγή συνδέεται υποχρεωτικά με ένα και μόνο ένα υπόλοιπο κρυπτονομίσματος.
    \item Ένα υπόλοιπο \selectlanguage{english}crypto\selectlanguage{greek} δεν μπορεί να υπάρχει 
    χωρίς αντίστοιχο πορτοφόλι.
    \item Κάθε ειδοποίηση τιμών αφορά αποκλειστικά ένα κρυπτονόμισμα.
    \item Δεν επιτρέπονται μεταφορές χρημάτων χωρίς συσχέτιση με έγκυρο τραπεζικό λογαριασμό.
    \item Οι τιμές \selectlanguage{english}candlestick\selectlanguage{greek} καταγράφονται ανά σταθερά 
    χρονικά διαστήματα.
    \item Κάθε \selectlanguage{english}candlestick\selectlanguage{greek} αντιστοιχεί σε ένα μόνο κρυπτονόμισμα.
\end{itemize}

\subsection{Καθορισμός Οντοτήτων}

% =============================
% USER
% =============================

\begin{tabular}{|p{4cm}|p{10cm}|}
\hline
\textbf{Όνομα Οντότητας} & \selectlanguage{english}User\selectlanguage{greek} \\ \hline
\textbf{Περιγραφή} & Αντιπροσωπεύει έναν εγγεγραμμένο χρήστη του συστήματος. \\ \hline
\textbf{Ιδιότητες} & Ισχυρή Οντότητα \\ \hline
\multirow{2}{4cm}{\textbf{Γνωρίσματα}} &
\selectlanguage{english}\underline{user\_id}\selectlanguage{greek}\\
\cline{2-2}
& \selectlanguage{english}\underline{email}\selectlanguage{greek}\\
\cline{2-2}
& \selectlanguage{english}name\selectlanguage{greek}\\
\cline{2-2}
& \selectlanguage{english}registered\_at\selectlanguage{greek}\\
\cline{2-2}
& \selectlanguage{english}phone\selectlanguage{greek}\\
\cline{2-2}
& \selectlanguage{english}password\_hash\selectlanguage{greek}\\
\hline
\end{tabular}

\vspace{1cm}

% =============================
% BANKACCOUNT
% =============================

\begin{tabular}{|p{4cm}|p{10cm}|}
\hline
\textbf{Όνομα Οντότητας} & \selectlanguage{english}BankAccount\selectlanguage{greek} \\ \hline
\textbf{Περιγραφή} & Αντιπροσωπεύει έναν τραπεζικό λογαριασμό που συνδέεται με έναν χρήστη. \\ \hline
\textbf{Ιδιότητες} & Ισχυρή Οντότητα \\ \hline
\multirow{2}{4cm}{\textbf{Γνωρίσματα}} &
\selectlanguage{english}\underline{account\_id}\selectlanguage{greek} \\
\cline{2-2}
& \selectlanguage{english}encrypted\_iban\selectlanguage{greek} \\
\cline{2-2}
& \selectlanguage{english}holder\_name\selectlanguage{greek} \\
\cline{2-2}
& \selectlanguage{english}bank\_name\selectlanguage{greek} \\
\cline{2-2}
& \selectlanguage{english}provider\_token\selectlanguage{greek} \\
\hline
\end{tabular}

\vspace{1cm}

% =============================
% WALLET
% =============================

\begin{tabular}{|p{4cm}|p{10cm}|}
\hline
\textbf{Όνομα Οντότητας} & \selectlanguage{english}Wallet\selectlanguage{greek} \\ \hline
\textbf{Περιγραφή} & Ένα ψηφιακό πορτοφόλι που ανήκει σε έναν ή περισσότερους χρήστες. \\ \hline
\textbf{Ιδιότητες} & Ισχυρή Οντότητα \\ \hline
\multirow{2}{4cm}{\textbf{Γνωρίσματα}} &
\selectlanguage{english}\underline{wallet\_id}\selectlanguage{greek} \\
\cline{2-2}
& \selectlanguage{english}cash\_balance\selectlanguage{greek} \\
\cline{2-2}
& \selectlanguage{english}name\selectlanguage{greek} \\
\hline
\end{tabular}

\vspace{1cm}

% =============================
% CRYPTOCURRENCY
% =============================

\begin{tabular}{|p{4cm}|p{10cm}|}
\hline
\textbf{Όνομα Οντότητας} & \selectlanguage{english}Cryptocurrency\selectlanguage{greek} \\ \hline
\textbf{Περιγραφή} & Αντιπροσωπεύει ένα κρυπτονόμισμα που παρακολουθεί το σύστημα. \\ \hline
\textbf{Ιδιότητες} & Ισχυρή Οντότητα \\ \hline
\multirow{2}{4cm}{\textbf{Γνωρίσματα}} &
\selectlanguage{english}\underline{ticker\_symbol}\selectlanguage{greek} \\
\cline{2-2}
& \selectlanguage{english}name\selectlanguage{greek} \\
\cline{2-2}
& \selectlanguage{english}logo\_image\selectlanguage{greek} \\
\hline
\end{tabular}

\vspace{1cm}

% =============================
% CANDLESTICK
% =============================

\begin{tabular}{|p{4cm}|p{10cm}|}
\hline
\textbf{Όνομα Οντότητας} & \selectlanguage{english}Candlestick\selectlanguage{greek} \\ \hline
\textbf{Περιγραφή} & Αναπαριστά \selectlanguage{english}candlestick\selectlanguage{greek} δεδομένα 
τιμών ενός κρυπτονομίσματος. \\ \hline
\textbf{Ιδιότητες} & Ασθενής Οντότητα (για τον πλήρη ορισμό της χρειάζεται
πέρα από το \selectlanguage{english}\dashuline{candle\_id}\selectlanguage{greek} και το 
\selectlanguage{english}\underline{ticker\_symbol}\selectlanguage{greek})\\ \hline
\multirow{2}{4cm}{\textbf{Γνωρίσματα}}
& \selectlanguage{english}\underline{ticker\_symbol}\selectlanguage{greek} \\
\cline{2-2}
& \selectlanguage{english}\dashuline{candle\_id}\selectlanguage{greek} \\
\cline{2-2}
& \selectlanguage{english}open\_price\selectlanguage{greek} \\
\cline{2-2}
& \selectlanguage{english}close\_price\selectlanguage{greek} \\
\cline{2-2}
& \selectlanguage{english}low\_price\selectlanguage{greek} \\
\cline{2-2}
& \selectlanguage{english}high\_price\selectlanguage{greek} \\
\cline{2-2}
& \selectlanguage{english}volume\selectlanguage{greek} \\
\cline{2-2}
& \selectlanguage{english}interval\selectlanguage{greek} \\
\hline
\end{tabular}

\vspace{1cm}

% =============================
% INVITATION
% =============================

\begin{tabular}{|p{4cm}|p{10cm}|}
\hline
\textbf{Όνομα Οντότητας} & \selectlanguage{english}Invitation\selectlanguage{greek} \\ \hline
\textbf{Περιγραφή} & Πρόσκληση που αποστέλλεται από έναν χρήστη σε έναν άλλον, 
ώστε να αποκτήσει δικαίωμα συνδικαιούχου σε ένα πορτοφόλι. \\ \hline
\textbf{Ιδιότητες} & Ισχυρή Οντότητα \\ \hline
\multirow{2}{4cm}{\textbf{Γνωρίσματα}} &
\selectlanguage{english}\underline{invitation\_id}\selectlanguage{greek} \\
\cline{2-2}
& \selectlanguage{english}status\selectlanguage{greek} \\
\cline{2-2}
& \selectlanguage{english}expires\_at\selectlanguage{greek} \\
\cline{2-2}
& \selectlanguage{english}created\_at\selectlanguage{greek} \\
\hline
\end{tabular}

\vspace{1cm}

% =============================
% CRYPTOHISTORY
% =============================

\begin{tabular}{|p{4cm}|p{10cm}|}
\hline
\textbf{Όνομα Οντότητας} & \selectlanguage{english}CryptoHistory\selectlanguage{greek} \\ \hline
\textbf{Περιγραφή} & Ιστορικά δεδομένα τιμών και όγκων ενός κρυπτονομίσματος. \\ \hline
\textbf{Ιδιότητες} & Ασθενής Οντότητα (για τον πλήρη ορισμό της χρειάζεται πέρα από το \selectlanguage{english}\dashuline{history\_id}\selectlanguage{greek} και το 
\selectlanguage{english}\underline{ticker\_symbol}\selectlanguage{greek})\\ \hline
\multirow{4}{4cm}{\textbf{Γνωρίσματα}} &
\selectlanguage{english}\underline{ticker\_symbol}\selectlanguage{greek} \\
\cline{2-2}
& \selectlanguage{english}\dashuline{history\_id}\selectlanguage{greek} \\
\cline{2-2}
& \selectlanguage{english}current\_price\selectlanguage{greek}\\
\cline{2-2}
& \selectlanguage{english}circulating\_supply\selectlanguage{greek}\\
\cline{2-2}
& \selectlanguage{english}timestamp\selectlanguage{greek}\\
\hline
\end{tabular}

\vspace{1cm}

% =============================
% CRYPTOBALANCE
% =============================

\begin{tabular}{|p{4cm}|p{10cm}|}
\hline
\textbf{Όνομα Οντότητας} & \selectlanguage{english}CryptoBalance\selectlanguage{greek} \\ \hline
\textbf{Περιγραφή} & Αντιπροσωπεύει το συνολικό ποσό κρυπτονομισμάτων που κατέχει ένα συγκεκριμένο \selectlanguage{english}wallet. \\ \hline
\textbf{Ιδιότητες} & Ασθενής Οντότητα (για τον πλήρη ορισμό της χρειάζεται πέρα από το \selectlanguage{english}\dashuline{balance\_id}\selectlanguage{greek} και το 
\selectlanguage{english}\underline{wallet\_id}\selectlanguage{greek})\\ \hline
\multirow{2}{4cm}{\textbf{Γνωρίσματα}} &
\selectlanguage{english}\underline{wallet\_id}\selectlanguage{greek} \\
\cline{2-2}
& \selectlanguage{english}\dashuline{balance\_id}\selectlanguage{greek} \\
\cline{2-2}
& \selectlanguage{english}crypto\_amount\selectlanguage{greek} \\
\cline{2-2}
& \selectlanguage{english}crypto\_balance\selectlanguage{greek} \\
\hline
\end{tabular}

\newpage
\subsection{Καθορισμός Συσχετίσεων}

\begin{tabular}{|p{4cm}|p{10cm}|}
\hline
\textbf{Όνομα Συσχέτισης} & \selectlanguage{english}User has Wallet\selectlanguage{greek} \\ \hline
\textbf{Περιγραφή} & Κάθε χρήστης διαθέτει ένα ή περισσότερα πορτοφόλια. Κάθε πορτοφόλι αντιστοιχεί σε έναν
ή περισσότερους χρήστες. \\ \hline
\textbf{Ιδιότητες} & \selectlanguage{english}Has-A\selectlanguage{greek} \\ \hline
\textbf{Λόγος Πληθικότητας} & M:N \\ \hline
\multirow{2}{4cm}{\textbf{Συμμετοχή}} 
& Ολική συμμετοχή του \selectlanguage{english}Wallet\selectlanguage{greek} \\
\cline{2-2}
& Μερική συμμετοχή του \selectlanguage{english}User\selectlanguage{greek} \\
\hline
\textbf{Γνωρίσματα} & \selectlanguage{english}role\selectlanguage{greek} \\ \hline
\end{tabular}

\vspace{1cm}

\begin{tabular}{|p{4cm}|p{10cm}|}
\hline
\textbf{Όνομα Συσχέτισης} & \selectlanguage{english}User owns BankAccount\selectlanguage{greek} \\ \hline
\textbf{Περιγραφή} & Ένας χρήστης μπορεί να διαθέτει έναν ή περισσότερους τραπεζικούς λογαριασμούς. 
Κάθε τραπεζικός λογαριασμός ανήκει σε έναν μόνο χρήστη. \\ \hline
\textbf{Ιδιότητες} & \selectlanguage{english}Has-A\selectlanguage{greek} \\ \hline
\textbf{Λόγος Πληθικότητας} & 1:N \\ \hline
\multirow{2}{4cm}{\textbf{Συμμετοχή}}
& Ολική συμμετοχή του \selectlanguage{english}BankAccount\selectlanguage{greek} \\ \cline{2-2}
& Μερική συμμετοχή του \selectlanguage{english}User\selectlanguage{greek} \\ \hline
\textbf{Γνωρίσματα} & - \\ \hline
\end{tabular}

\vspace{1cm}

\begin{tabular}{|p{4cm}|p{10cm}|}
\hline
\textbf{Όνομα Συσχέτισης} & \selectlanguage{english}User receives Invitation\selectlanguage{greek} \\ \hline
\textbf{Περιγραφή} & Ένας χρήστης μπορεί να λάβει πολλές προσκλήσεις. Κάθε πρόσκληση στέλνεται σε 
έναν μόνο χρήστη. \\ \hline
\textbf{Ιδιότητες} & \selectlanguage{english}Has-A\selectlanguage{greek} \\ \hline
\textbf{Λόγος Πληθικότητας} & 1:N \\ \hline
\multirow{2}{4cm}{\textbf{Συμμετοχή}}
& Ολική συμμετοχή του \selectlanguage{english}Invitation\selectlanguage{greek} \\ \cline{2-2}
& Μερική συμμετοχή του \selectlanguage{english}User\selectlanguage{greek} \\ \hline
\textbf{Γνωρίσματα} & - \\ \hline
\end{tabular}

\vspace{1cm}

\begin{tabular}{|p{4cm}|p{10cm}|}
\hline
\textbf{Όνομα Συσχέτισης} & \selectlanguage{english}User sends Invitation\selectlanguage{greek} \\ \hline
\textbf{Περιγραφή} & Ένας χρήστης μπορεί να στείλει πολλές προσκλήσεις. Κάθε πρόσκληση στέλνεται από 
έναν μόνο χρήστη. \\ \hline
\textbf{Ιδιότητες} & \selectlanguage{english}Has-A\selectlanguage{greek} \\ \hline
\textbf{Λόγος Πληθικότητας} & 1:N \\ \hline
\multirow{2}{4cm}{\textbf{Συμμετοχή}}
& Ολική συμμετοχή του \selectlanguage{english}Invitation\selectlanguage{greek} \\ \cline{2-2}
& Μερική συμμετοχή του \selectlanguage{english}User\selectlanguage{greek} \\ \hline
\textbf{Γνωρίσματα} & - \\ \hline
\end{tabular}

\vspace{1cm}

\begin{tabular}{|p{4cm}|p{10cm}|}
\hline
\textbf{Όνομα Συσχέτισης} & \selectlanguage{english}User gets\_alert\_from Cryptocurrency \\ \hline
\textbf{Περιγραφή} & Κάθε χρήστης μπορεί να δέχεται  \selectlanguage{english}alerts\selectlanguage{greek} από πολλα 

\selectlanguage{english}cryptocurrencys\selectlanguage{greek}. Κάθε \selectlanguage{english}cryptocurrency\selectlanguage{greek} στέλνει \selectlanguage{english}alert\selectlanguage{greek} σε έναν ή περισσότερους χρήστες.\\ \hline
\textbf{Ιδιότητες} & \selectlanguage{english}Has-A\selectlanguage{greek} \\ \hline
\textbf{Λόγος Πληθικότητας} & M:N \\ \hline
\multirow{2}{4cm}{\textbf{Συμμετοχή}}
& Ολική συμμετοχή του \selectlanguage{english}Cryptocurrency\selectlanguage{greek} \\ \cline{2-2}
& Ολική συμμετοχή του \selectlanguage{english}User\selectlanguage{greek} \\ \hline
\textbf{Γνωρίσματα} & \selectlanguage{english}threshold\_direction, threshold\_price, is\_active, 

triggered\_at \\ \hline
\end{tabular}

\vspace{1cm}

\begin{tabular}{|p{4cm}|p{10cm}|}
\hline
\textbf{Όνομα Συσχέτισης} & \selectlanguage{english}User transacts CryptoBalance\selectlanguage{greek} \\ \hline
\textbf{Περιγραφή} & Κάθε χρήστης μπορεί να πραγματοποιεί πολλές συναλλαγές. Κάθε συναλλαγή δημιουργείται από έναν ή παραπάνω χρήστες (συνδικαιούχοι ενός πορτοφολιού). \\ \hline
\textbf{Ιδιότητες} & \selectlanguage{english}Has-A\selectlanguage{greek} \\ \hline
\textbf{Λόγος Πληθικότητας} & M:N \\ \hline
\multirow{2}{4cm}{\textbf{Συμμετοχή}} 
& Μερική συμμετοχή του \selectlanguage{english}User\selectlanguage{greek} \\
\cline{2-2}
& Ολική συμμετοχή του \selectlanguage{english}CryptoBlance\selectlanguage{greek} \\
\hline
\textbf{Γνωρίσματα} & \selectlanguage{english}amount, price\_at\_timestamp, timestamp, 

transaction\_type, transaction\_hash \\ \hline
\end{tabular}

\vspace{1cm}

\begin{tabular}{|p{4cm}|p{10cm}|}
\hline
\textbf{Όνομα Συσχέτισης} & \selectlanguage{english}BankAccount transfers\_to Wallet\selectlanguage{greek} \\ \hline
\textbf{Περιγραφή} & Κάθε τραπεζικός λογαριασμός μεταφέρει χρήματα σε κάποιο πορτοφόλι. Κάθε πορτοφόλι δέχεται μεταφορές από κάποιον τραπεζικό λογαριασμό.\\ \hline
\textbf{Ιδιότητες} & \selectlanguage{english}Has-A\selectlanguage{greek}, Προσδιορίζουσα \\ \hline
\textbf{Λόγος Πληθικότητας} & M:N \\ \hline
\multirow{2}{4cm}{\textbf{Συμμετοχή}}
& Ολική συμμετοχή του \selectlanguage{english}Candlestick\selectlanguage{greek} \\ \cline{2-2}
& Ολική συμμετοχή του \selectlanguage{english}Cryptocurrency\selectlanguage{greek} \\ \hline
\textbf{Γνωρίσματα} & \selectlanguage{english}amount, timestamp, direction \\ \hline
\end{tabular}

\vspace{1cm}

\begin{tabular}{|p{4cm}|p{10cm}|}
\hline
\textbf{Όνομα Συσχέτισης} & \selectlanguage{english}Invitation refers\_to Wallet\selectlanguage{greek} \\ \hline
\textbf{Περιγραφή} & Κάθε πρόσκληση αφορά ένα συγκεκριμένο πορτοφόλι. Ένα πορτοφόλι μπορεί να αποτελεί αντικείμενο πολλών προσκλήσεων. \\ \hline
\textbf{Ιδιότητες} & \selectlanguage{english}Has-A\selectlanguage{greek} \\ \hline
\textbf{Λόγος Πληθικότητας} & Ν:1 \\ \hline
\multirow{2}{4cm}{\textbf{Συμμετοχή}} 
& Ολική συμμετοχή του \selectlanguage{english}Invitation\selectlanguage{greek} \\
\cline{2-2}
& Μερική συμμετοχή του \selectlanguage{english}Wallet\selectlanguage{greek} \\
\hline
\textbf{Γνωρίσματα} & - \\ \hline
\end{tabular}

\vspace{1cm}

\begin{tabular}{|p{4cm}|p{10cm}|}
\hline
\textbf{Όνομα Συσχέτισης} & \selectlanguage{english}Wallet has CryptoBalance\selectlanguage{greek} \\ \hline
\textbf{Περιγραφή} & Ένα πορτοφόλι μπορεί να διαθέτει υπόλοιπα για πολλά κρυπτονομίσματα. Κάθε 
\selectlanguage{english}cryptobalance\selectlanguage{greek} ανήκει σε ένα μόνο πορτοφόλι. \\ \hline
\textbf{Ιδιότητες} & \selectlanguage{english}Has-A\selectlanguage{greek} \\ \hline
\textbf{Λόγος Πληθικότητας} & 1:N \\ \hline
\multirow{2}{4cm}{\textbf{Συμμετοχή}} 
& Ολική συμμετοχή του \selectlanguage{english}CryptoBalance\selectlanguage{greek} \\
\cline{2-2}
& Μερική συμμετοχή του \selectlanguage{english}Wallet\selectlanguage{greek} \\
\hline
\textbf{Γνωρίσματα} & - \\ \hline
\end{tabular}

\vspace{1cm}

\begin{tabular}{|p{4cm}|p{10cm}|}
\hline
\textbf{Όνομα Συσχέτισης} & \selectlanguage{english}CryptoBalance involves Cryptocurrency
\selectlanguage{greek} \\ \hline
\textbf{Περιγραφή} & Κάθε υπόλοιπο κρυπτονομίσματος αντιστοιχεί σε ένα συγκεκριμένο κρυπτονόμισμα. 
Ένα κρυπτονόμισμα μπορεί να σχετίζεται με πολλά υπόλοιπα κρυπτονομισμάτων. \\ \hline
\textbf{Ιδιότητες} & \selectlanguage{english}Has-A\selectlanguage{greek} \\ \hline
\textbf{Λόγος Πληθικότητας} & 1:N \\ \hline
\multirow{2}{4cm}{\textbf{Συμμετοχή}} 
& Ολική συμμετοχή του \selectlanguage{english}CryptoBalance\selectlanguage{greek} \\
\cline{2-2}
& Μερική συμμετοχή του \selectlanguage{english}Cryptocurrency\selectlanguage{greek} \\
\hline
\textbf{Γνωρίσματα} & - \\ \hline
\end{tabular}

\vspace{1cm}

\begin{tabular}{|p{4cm}|p{10cm}|}
\hline
\textbf{Όνομα Συσχέτισης} & \selectlanguage{english}Cryptocurrency has Candlestick\selectlanguage{greek} \\ \hline
\textbf{Περιγραφή} & Κάθε κρυπτονόμισμα έχει πολλά \selectlanguage{english}candlesticks\selectlanguage{greek}. 
Κάθε \selectlanguage{english}candlestick\selectlanguage{greek} αντιστοιχεί σε ένα μόνο κρυπτονόμισμα. \\ \hline
\textbf{Ιδιότητες} & \selectlanguage{english}Has-A\selectlanguage{greek}, Προσδιορίζουσα \\ \hline
\textbf{Λόγος Πληθικότητας} & 1:N \\ \hline
\multirow{2}{4cm}{\textbf{Συμμετοχή}}
& Ολική συμμετοχή του \selectlanguage{english}Candlestick\selectlanguage{greek} \\ \cline{2-2}
& Μερική συμμετοχή του \selectlanguage{english}Cryptocurrency\selectlanguage{greek} \\ \hline
\textbf{Γνωρίσματα} & - \\ \hline
\end{tabular}

\vspace{1cm}

\begin{tabular}{|p{4cm}|p{10cm}|}
\hline
\textbf{Όνομα Συσχέτισης} & \selectlanguage{english}Cryptocurrency has CryptoHistory\selectlanguage{greek} \\ \hline
\textbf{Περιγραφή} & Ένα κρυπτονόμισμα έχει πολλά ιστορικά σημεία τιμών. Κάθε ιστορικό σημείο αντιστοιχεί σε ένα κρυπτονόμισμα. \\ \hline
\textbf{Ιδιότητες} & \selectlanguage{english}Has-A\selectlanguage{greek} \\ \hline
\textbf{Λόγος Πληθικότητας} & 1:N \\ \hline
\multirow{2}{4cm}{\textbf{Συμμετοχή}} 
& Ολική συμμετοχή του \selectlanguage{english}CryptoHistory\selectlanguage{greek} \\
\cline{2-2}
& Μερική συμμετοχή του \selectlanguage{english}Cryptocurrency\selectlanguage{greek} \\
\hline
\textbf{Γνωρίσματα} & - \\ \hline
\end{tabular}

\newpage
% % !!!!!!!!!!!!!  The ER diagram is in the second .tex file  !!!!!!!!!!!!!

% \subsection{Διάγραμμα Οντοτήτων/Συσχετίσεων}

% \begin{figure}[ht]
%     \centering
%     \includegraphics[width=1\linewidth]{db_deliverable_1.png}
%     \label{fig:ER}
% \end{figure}



\end{document}
